% Integrated discussion replacing the article's overlapping Discussion and Conclusion.
\section{What retrieval contributes to learning from sparse evidence}
\label{sec:ch4:thesis-discussion}

\subsection{Selective compensation at inference time}

The main contribution of this chapter is a controlled change in how linguistic evidence is used. For a given trained base model and evaluation set, adding $k$NN retrieval changes the probabilities assigned to the two members of a syntactic minimal pair. Low-frequency items improve in all three phenomenon groups and both model settings in the representative configuration. For example, subject--verb agreement accuracy in the smaller model rises from 0.60 to 0.76 for low-frequency items, while high-frequency accuracy changes from 0.91 to 0.90 (Table~\ref{tab:frequency_effects}). The frequency gap therefore narrows primarily through improved low-frequency performance, although small declines in some high-frequency conditions also contribute to gap reduction.

Relative clauses show the largest absolute low-frequency gains in the reported table: from 0.48 to 0.65 in the smaller model and from 0.41 to 0.70 in GPT-2 XL. This pattern identifies relative-clause items as especially responsive to the tested intervention. It does not by itself establish a general law relating syntactic complexity to retrieval benefit. The phenomenon groups differ in baseline accuracy, lexical composition, and benchmark sources, and lower baselines leave more room for improvement. The stronger conclusion is that retrieval can supply useful evidence in conditions where the base model's decisions are least reliable. Low-frequency accuracy nevertheless remains below high-frequency accuracy after augmentation.

This result concerns the contribution of retrieval at inference time. Interpolation can improve a decision even when the base model's parameters remain unchanged. The experiment therefore distinguishes access to useful evidence from its prior incorporation into the parametric predictor, but does not measure subsequent consolidation or durable learning after the datastore is removed. Likewise, lower perplexity and higher minimal-pair accuracy are complementary outcomes: improved average prediction does not automatically explain improved discrimination of a particular grammatical contrast.

\subsection{The usefulness of memory depends on its contents and representation}

The ablations show that the effect depends on how stored experience is represented and selected. Averaged contextual representations and the GloVe and fastText alternatives generally fail to reproduce the benefits of full contextual retrieval (Figure~\ref{fig:vector}). In the structural-content analysis, retrieval sets containing the target phenomenon are more useful for low-frequency items than sets without it, and similarity ranking does not fully compensate for the absence of those examples (Figure~\ref{fig:datastore}). A useful memory must therefore make task-relevant information accessible, rather than merely return lexically related material.

These comparisons do not cleanly separate syntax from semantics. Averaging contextual representations removes positional distinctions in the pooling operation but may retain information about structure encoded before pooling. Static lexical embeddings differ from contextual embeddings in several properties, including their training sources and similarity geometry. Their poorer performance cannot be attributed uniquely to the removal of syntax. Moreover, the structural-match indicator in Appendix~\ref{app:struc_simi} records whether a retrieved sequence contains the same broad phenomenon, using parser-based criteria. It is not a measure of identity between complete syntactic analyses or between the particular dependencies in two sentences. The results support the usefulness of these operationally defined structural cues; a stronger decomposition would require closer control of lexical content, construction, and representational quality.

Configuration matters as well. In the main sequence-level comparison, 16 neighbours outperform both 1 and 1,024 neighbours on the six low-frequency conditions. The additional analyses support an intermediate setting as a useful choice among the tested values, rather than a universal optimum. Context-size effects differ across phenomena, and the advantage of sequence over smaller retrieval units is modest. Deterioration with more neighbours is compatible with interference from less relevant instances, but the experiment does not isolate interference as the mechanism. Changes in the resulting token distribution and in its interaction with the fixed interpolation weight are alternative contributors. Adaptive retrieval is consequently a motivated future intervention, not a mechanism already demonstrated by this study.

\subsection{Cognitive and methodological implications}

For cognitive modelling, the study supplies a functional demonstration of complementarity. A system that combines parametric prediction with direct access to instances can make more accurate low-frequency syntactic decisions than its parametric component alone. This is consistent with the broad motivation for complementary learning systems, while leaving open whether children use a comparable retrieval process during the same tasks. The experiment contains no matched child retrieval manipulation, no neural measurement, and no developmental trajectory of memory formation. Hippocampal storage, replay, selective encoding, forgetting, and consolidation are therefore candidate mechanisms for further modelling, rather than findings of the present comparison.

For AI methods, the contribution is more specific than an overall gain from retrieval augmentation. Frequency-stratified evaluation shows where an apparently successful configuration helps, where it leaves a residual gap, and where it can reduce already strong performance. The representation and retrieval-content ablations then identify conditions under which that compensation is lost. This combination turns a general retrieval method into a diagnostic instrument for examining the reliability of syntactic predictions across lexical frequencies. It also motivates evaluating retrieval strategies on both aggregate accuracy and the distribution of gains across items.

The developmental interpretation depends on a complete account of the system's information sources. The smaller base model was trained for ten epochs over a 50.03M-word corpus whose largest component is subtitles; this is bounded textual training, not a record of one child's successive encounters. The full-retrieval method is described using the base GPT-2 model's contextual states, while GloVe and fastText introduce different learned representations in the ablations. The datastore's corpus, construction split, size, and overlap with training and evaluation need to be reported explicitly. Only a datastore restricted to previously available material would isolate renewed access to that material; a store containing additional text would also change the evidence available to the system. This distinction is especially consequential for claims about learning from limited experience.
% AUTHOR CHECK: Verify datastore provenance, frequency-count corpus (especially for GPT-2 XL), tokenizer provenance, and external embedding checkpoints; do not substitute an undocumented DPR pipeline.

Finally, the evaluation measures probabilistic preference on selected written minimal pairs. BLiMP and the other sources provide controlled contrasts, but the present subsets do not establish mastery of all English syntax, comprehension in context, spontaneous grammatical production, or a child's developmental course. Frequency bands are defined by the mean frequency of content words in an item, so they should not be equated with an experimentally manipulated number of encounters with a single target word. Comparisons across the two model settings are further limited by simultaneous differences in model size, corpus composition, and training history.

\subsection{Targeted next steps}

Three interventions follow from these results. First, selective storage or weighting could prioritize rare or informative instances rather than treating similarity as the only retrieval criterion. Error-driven learning and selective data weighting motivate this comparison \citep{chang2006dual,rescorla1972theory,toneva2018empirical,xie2023doremi}. A useful test would keep the datastore's total size fixed while varying which instances are retained, then assess both low-frequency improvement and any loss on frequent items.

Second, retrieval could be constrained by the target construction or use alternative contextual representations. Comparing the present layer-normalized keys with other points in the residual stream, while controlling the datastore and distance scaling, would test whether a persistent gap reflects selection of the wrong instances or inadequate information in the selected representation. Structural filtering would need matched lexical controls so that an apparent structural gain is not explained by a different distribution of words.

Third, an adaptive rule could choose context scope and neighbour count for the current prediction. The present configuration effects motivate such a rule, but do not establish its benefit. Comparing a rule selected on held-out development material with fixed configurations would test whether adaptation improves the distribution of gains without selecting settings on the final syntactic test.

Together with Chapter~\ref{chp2_chap}, these results make frequency a recurring question about the availability and use of linguistic evidence. The preceding chapter examined whether words appear reliably in generated output; this chapter shows that stored instances can selectively improve another expression of linguistic knowledge. A direct bridge would test retrieval on the same expressive-production measure, or evaluate one learner on both production and minimal pairs while controlling its information sources. The following chapter examines a different resource, caregiver interaction and feedback. These studies extend the thesis by testing distinct ways in which experience can influence model behaviour, without treating them as successive stages of one developing system.
